\documentclass[11pt]{article}
\usepackage{amssymb, latexsym, amsmath, verbatim, multicol}

\setlength{\oddsidemargin}{-.25in} \setlength{\evensidemargin}{-.25in} \setlength{\textwidth}{6.8in}
\setlength{\textheight}{9.3in} \setlength{\topmargin}{-.8in} \setlength\parindent{0pt}      

\newcommand{\midtermOrFinal}{Midterm 2}
\newcommand{\courseNumber}{271}
\newcommand{\courseName}{Mathematical Reasoning and Proof}
% \newcommand{\courseName}{Linear algebra}
\newcommand{\season}{Spring}
\newcommand{\thisYear}{2024}


%%%%%%%%%%%%%%%%%%%%%%%%%%%%%%%%%%%%%%%%%%%%%%%%%%%%%%%%%%%%%%%%%%%%%%%%%%%%%%%%%%%%%%%%%%%%%%%%%%%%%%%%%
% question counter
%%%%%%%%%%%%%%%%%%%%%%%%%%%%%%%%%%%%%%%%%%%%%%%%%%%%%%%%%%%%%%%%%%%%%%%%%%%%%%%%%%%%%%%%%%%%%%%%%%%%%%%%% 
\newcommand{\class}[1]{\left< #1 \right>}

\newtoks\scores
\newcounter{sumscore} \newcounter{quesno} \newcounter{inquesno} \newcounter{partno} \setcounter{inquesno}{1} \newcounter{altquesno}

\newcommand{\ques}[1]{
    \addtocounter{sumscore}{#1}
    \stepcounter{quesno}
    \expandafter\scores\expandafter{\the\scores
    \theinquesno\stepcounter{inquesno}&#1&\\ \hline}
    \par\noindent\thequesno. (#1 points)
}
\newcommand{\quest}[1]{
        \addtocounter{sumscore}{#1}
        \stepcounter{quesno}
        \expandafterscores\expandafter{\the\scores
        \theinquesno\stepcounter{inquesno}&#1&\\ \hline}
        \par\noindent\thequesno.(#1points)
}
%%%%%%%%%%%%%%%%%%%%%%%%%%%%%%%%%%%%%%%%%%%%%%%%%%%%%%%%%%%%%%%%%%%%%%%%%%%%%%%%%%%%%%%%%%%%%%%%%%%%%%%%% 

\newcommand{\N}{{\mathbb N}}
\newcommand{\Z}{{\mathbb Z}}
\newcommand{\C}{{\mathbb C}}
\newcommand{\R}{{\mathbb R}}
\newcommand{\Q}{{\mathbb Q}}
\DeclareMathOperator{\curl}{curl}
\newcommand{\Span}{{\rm Span}}

%%%%%%%%%%%%%%%%%%%%%%%%%%%%%%%%%%%%%%%%%%%%%%%%%%%%%%%%%%%%%%%%%%%%%%%%%%%%%%%%%%%%%%%%%%%%%%%%%%%%%%%%%

\begin{document}

\begin{multicols}{2}
  \scriptsize{\noindent
       {\sc Instructor: David Zureick-Brown}\\ 
        {\season } {\thisYear} 
        \columnbreak
        \begin{flushright}
        Amherst College\\
        Dept.~of Mathematics and Statistics\\
        [20pt]
        \end{flushright}}
\end{multicols}

%%%%%%%%%%%%%%%%%%%%%%%%%%%%%%%%%%%%%%%%%%%%%%%%%%%%%%%%%%%%%%%%%%%%%%%%%%%%%%%%%%%%%%%%%%%%%%%%%%%%%%%%% 

\vskip-.18in
\noindent\rule{6.8in}{1pt}
\begin{center}
\emph{{\bf         Math {\courseNumber } $\sim$ {\courseName }:  {\midtermOrFinal }}}\\[25pt]
\end{center}
\vskip-.28in

\noindent\rule{6.8in}{1pt}
\vskip.55in


%%%%%%%%%%%%%%%%%%%%%%%%%%%%%%%%%%%%%%%%%%%%%%%%%%%%%%%%%%%%%%%%%%%%%%%%%%%%%%%%%%%%%%%%%%%%%%%%%%%%%%%%% 
\textbf{Name} (print clearly as it appears on Gradescope)\textbf{:} .....................................................................................
	
\vskip.25in
\begin{itemize}
\item This is a closed-book, 50-minute examination. 

\item No books, notes, calculators, cell phones, communication devices of any sort, or other aids are permitted.

\item I will be in the hallway, and will check in 20 and 40 minutes for questions.  
  
\item You may leave when you are finished, in which case please hand me the exam in the hallway.

\item You may leave to use the bathroom.
  
\item You may use the backs of pages for additional work space. If you use any scratch paper as part of your solution, please let me know after the exam.

\item Be sure to justify your answer to each of the problems!

\item For proofs, please make sure that your answers are in complete sentences!

\item {\Large \it Good luck!}  \vspace{1cm}
\end{itemize}


\newpage
%%%%%%%%%%%%%%%%%%%%%%%%%%%%%%%%%%%%%%%%%%%%%%%%%%%%%%%%%%%%%%%%%%%%%%%%%%%%%%%%%%%%%%%%%%%%%%%%%%%%%%%%% 
\ques{15} Let $f \colon A \to B$ be a function.
 Define what it means for $f$ to be \emph{injective}.
% Define what it means for $f$ to be \emph{surjective}.

\vspace{2.5in}
 
Negate the definition of injective.\\

% For all $x \in \mathbb{R}$, if there exists $y \in \mathbb{R}$ such that $xy = 1$, then $x \neq 0$.

% There exists  $x \in \mathbb{R}$ such that if $x \neq 0$, then  for all  $y \in \mathbb{R}$,  $xy = 1$.


% \vspace{3cm}
% \ques{5} Give an example of a function that is surjective but not injective.

\vspace{3.5in}


% Let $f\colon A \to B$ be a function and let $X \subseteq A$. Define the \emph{image} $f(X)$ of $X$.
 % Let $f\colon A \to B$ be a function and let $W \subseteq B$. Define the \emph{preimage} $f^{-1}(W)$ of $W$.

Let $f\colon A \to B$ be a function and let $X \subseteq A$. We define the \emph{image} $f(X)$ of $X$ to be
\[
f(X) = \{\hspace{20in}
  \]

% Let $f\colon A \to B$ be a function and let $W \subseteq B$. We define the \emph{preimage} $f^{-1}(W)$ of $W$ to be 
% \[
% f^{-1}(W) = \{\hspace{20in}
%   \]


\newpage
%%%%%%%%%%%%%%%%%%%%%%%%%%%%%%%%%%%%%%%%%%%%%%%%%%%%%%%%%%%%%%%%%%%%%%%%%%%%%%%%%%%%%%%%%%%%%%%%%%%%%%%%%
\ques{15}   Label each of the following functions as injective, surjective, 
  bijective, or neither. (No proof is necessary.)\vspace{10pt}

  
  \begin{enumerate}


  \item \textbf{I}nj\hspace{5pt} \textbf{S}urj \, \textbf{B}ij \, \textbf{N}:
    \hspace{12pt}
$\mathbb{R} \to \mathbb{R};\, x \mapsto x(x^2-1)$.
\smallskip
  \item \textbf{I}nj\hspace{5pt} \textbf{S}urj \, \textbf{B}ij \, \textbf{N}:
    \hspace{12pt} 
$\mathbb{R}_{ > 0} \to \mathbb{R}$,  $ x \mapsto \ln \left(1/x\right)  .$
\smallskip
\item \textbf{I}nj\hspace{5pt} \textbf{S}urj \, \textbf{B}ij \, \textbf{N}:
    \hspace{12pt}
$P(\mathbb{Z}) \to P(\mathbb{Z});\, S \mapsto S \cap \{1\}$.

\vspace{22pt}



%   \item \textbf{I}nj\hspace{5pt} \textbf{S}urj \, \textbf{B}ij \, \textbf{N}:
%     \hspace{12pt}
% $P(\mathbb{Z}) \to P(\mathbb{Z});\, S \mapsto S \cap \{1\}$.
% \vspace{22pt}

\end{enumerate}

\ques{15} 
% Let $f \colon \mathbb{R}-\{0\} \to \mathbb{R}$ be the function
%  $$ x \mapsto {\frac{x - 1}{x + 1}}.$$
%  \begin{enumerate}
%  \item Prove that $f$ is injective.
%  \item Prove that $f$ is not surjective.
%  \end{enumerate}
Let $f \colon \mathbb{R} \to \mathbb{R}$ be the function
 $$ x \mapsto e^{(x - 1)(x+1)} .$$
 \begin{enumerate}
 \item Prove that $f$ is not injective.
 \item Prove that $f$ is not surjective.
 \end{enumerate}

 % Let $f \colon \mathbb{R}-\{-1\} \to \mathbb{R}$ be the function
 % $$ x \mapsto e^{(x^2 - 1)} .$$
 % \begin{enumerate}
 % \item Prove that $f$ is injective.
 % \item Prove that $f$ is not surjective.
 % \end{enumerate}



% $P(\mathbb{Z}) \to \mathbb{Z};\, S \mapsto |S| \text{ if } S \text{ is finite}, 0
% \text{ if } S \text{ is infinite}$.


% \vspace{2.5in}

% \newpage
% %%%%%%%%%%%%%%%%%%%%%%%%%%%%%%%%%%%%%%%%%%%%%%%%%%%%%%%%%%%%%%%%%%%%%%%%%%%%%%%%%%%%%%%%%%%%%%%%%%%%%%%%% 
% \ques{5} Circle whichever of the following statements are true. (No proof is necessary.)\\

%   \begin{enumerate}
%   % \item $\{0\} \subseteq P(\mathbb{R})$\vspace{0.3cm}
%   \item $\mathbb{Z} \subseteq P(\mathbb{R})$\vspace{0.3cm}
%   \item $\mathbb{Z} \in P(\mathbb{R})$\vspace{0.3cm}
%   % \item $\emptyset \in P(\mathbb{R})$\vspace{0.3cm}
%   % \item $\emptyset \subseteq P(\mathbb{R})$\vspace{0.3cm}
%   \item $\{\emptyset\} \in P(\mathbb{R})$ \vspace{0.3cm}
%   \item $\{\emptyset\} \subseteq P(\mathbb{R})$\vspace{0.3cm}
%    \item $\{\{1\},\{1,2\}\} \subseteq P(\mathbb{R})$\\
%   \end{enumerate}



% \newpage
% %%%%%%%%%%%%%%%%%%%%%%%%%%%%%%%%%%%%%%%%%%%%%%%%%%%%%%%%%%%%%%%%%%%%%%%%%%%%%%%%%%%%%%%%%%%%%%%%%%%%%%%%% 
% \ques{15} 
% Prove or disprove. (For disproofs, please also give a counterexample.)
% \begin{enumerate}
% \item For all sets $A,B,C$,  $A \cup (B \cap C) = (A\cup B) \cap C$.

% \item For all sets $A,B,C$, if $A-B \subseteq C$, then $\overline{C} \subseteq \overline{A} \cup B$.
% \end{enumerate}


\newpage
%%%%%%%%%%%%%%%%%%%%%%%%%%%%%%%%%%%%%%%%%%%%%%%%%%%%%%%%%%%%%%%%%%%%%%%%%%%%%%%%%%%%%%%%%%%%%%%%%%%%%%%%%
\ques{10}  Let $f\colon A \to B$ and $g \colon B \to C$ be functions.
Prove that  if $f$ and $g$ are injective, then $g\circ f$ is injective.
% Prove that  if $f$ and $g$ are surjective, then $g\circ f$ is surjective.

\newpage
%%%%%%%%%%%%%%%%%%%%%%%%%%%%%%%%%%%%%%%%%%%%%%%%%%%%%%%%%%%%%%%%%%%%%%%%%%%%%%%%%%%%%%%%%%%%%%%%%%%%%%%%%
\ques{10}  
State the converse of the previous problem and disprove it by giving a counterexample.



% \item   Let $f\colon A \to B$ and $g \colon B \to C$ be functions. Prove or
%   disprove each of the following:
%   \begin{enumerate}
%   \item If $f$ and $g$ are injections, then $gf$ is an injection.
%   \item If $f$ and $g$ are surjections, then $gf$ is a surjection.
%   \item If $f$ and $g$ are bijections, then $gf$ is a bijection.
%   \item If $gf$ is an injection, then $f$ and $g$ are injections.
%   \item If $gf$ is a surjection, then $f$ and $g$ are surjections.
%   \item If $gf$ is a bijection, then $f$ and $g$ are bijections.
%   \item If $gf$ is an injection, then $f$ is an injection.
%   \item If $gf$ is an injection, then $g$ is an injection.
%   \item If $gf$ is a surjection, then $f$ is a surjection.
%   \item If $gf$ is a surjection, then $g$ is a surjection.
%   \item If $gf$ is a bijection, then $f$ is a bijection.
%   \item If $gf$ is a bijection, then $g$ is a bijection.
%   \item If $gf$ is an injection and $g$ is invertible, then $f$ is an
%     injection.

%   \item If $gf$ is an injection and $f$ is surjective, then $g$ is an
%     injection.

%   \item If $gf$ is a surjection and $g$ is injective, then $f$ is an
%  surjective.

%   \end{enumerate}


\newpage
%%%%%%%%%%%%%%%%%%%%%%%%%%%%%%%%%%%%%%%%%%%%%%%%%%%%%%%%%%%%%%%%%%%%%%%%%%%%%%%%%%%%%%%%%%%%%%%%%%%%%%%%%
\ques{15}
Let $f\colon A \to B$ be a function and let $X,Y \subseteq B$. Prove that $X \subseteq Y \Rightarrow f^{-1}(X) \subseteq f^{-1}(Y)$.



\newpage
%%%%%%%%%%%%%%%%%%%%%%%%%%%%%%%%%%%%%%%%%%%%%%%%%%%%%%%%%%%%%%%%%%%%%%%%%%%%%%%%%%%%%%%%%%%%%%%%%%%%%%%%%
\ques{10}
Let $f\colon A \to B$ be a function and let $X,Y \subseteq B$. Prove that the statement ``$f^{-1}(X) \subseteq f^{-1}(Y) \Rightarrow X \subseteq Y$'' is false by giving a counterexample.
   
\newpage
%%%%%%%%%%%%%%%%%%%%%%%%%%%%%%%%%%%%%%%%%%%%%%%%%%%%%%%%%%%%%%%%%%%%%%%%%%%%%%%%%%%%%%%%%%%%%%%%%%%%%%%%%
\ques{10}
Let $f\colon A \to B$ be a function and let $X,Y \subseteq B$. Prove that if $f$ is surjective then $f^{-1}(X) \subseteq f^{-1}(Y) \Rightarrow X \subseteq Y$.



% \newpage
% %%%%%%%%%%%%%%%%%%%%%%%%%%%%%%%%%%%%%%%%%%%%%%%%%%%%%%%%%%%%%%%%%%%%%%%%%%%%%%%%%%%%%%%%%%%%%%%%%%%%%%%%% 
% \ques{15} Let $f \colon A \to B$ be a function and let $W \subset
% B$. Prove or disprove each of the following. (For disproofs, please also give a counterexample.)
% \begin{enumerate}
%   \item $W \subseteq f(f^{-1}(W))$.
%   \item $W \supseteq f(f^{-1}(W))$.
% \end{enumerate}

% \ques{20} Let $f \colon A \to B$ be a function and let $X,Y \subset
% A$. Prove or disprove each of the following. (For disproofs, please also give a counterexample.)
% \begin{enumerate}
% \item  $f(X - Y) \subseteq f(X) - f(Y)$.
% \item  $f(X - Y) \supseteq f(X) - f(Y)$.
% \end{enumerate}


% \ques{20} Let $f \colon A \to B$ be a function and let $W \subset
% A$. Prove or disprove each of the following. (For disproofs, please also give a counterexample.)
% \begin{enumerate}
%   \item $W \subseteq f^{-1}(f(W))$.
%   \item $W \supseteq f^{-1}(f(W))$.
% \end{enumerate}





% \newpage
% %%%%%%%%%%%%%%%%%%%%%%%%%%%%%%%%%%%%%%%%%%%%%%%%%%%%%%%%%%%%%%%%%%%%%%%%%%%%%%%%%%%%%%%%%%%%%%%%%%%%%%%%% 

% \ques{10}  
% Let $f\colon A \to B$ be a function and let $X,Y \subset
%   A$. Prove that if $f$ is injective then $f(X \cap Y) \supseteq f(X) \cap f(Y)$.

% \newpage
% \ques{15}  Which of the following binary operations are
%   commutative/associative, which have an identity element, and for
%   which does every element have an inverse? (Do not prove that your
%   answers are correct.)

%   \begin{enumerate}
%  % \item $\Z \times \Z \xrightarrow{+} \Z$.
%  %  \item $\R \times \R \xrightarrow{\cdot} \R$.
%  %  \item $\R \times \R \xrightarrow{/} \R$.
%   % \item $\R^* \times \R^* \xrightarrow{/} \R^*$ (where $\R^* = \R -    \{0\}$.)
%   \item $Fun(B,B) \times Fun(B,B) \xrightarrow{\circ} Fun(B,B)$.
%   % \item $\Z \times \Z \xrightarrow{\star} \Z$ (where $a \star b = a +
%   %   b + 1$).
%   \item $\Z \times \Z \xrightarrow{\star} \Z$ (where $a \star b = 2a
%      +b$).
%   % \item $P(A)\times P(A) \xrightarrow{\cap} P(A)$.
%   \item $P(A)\times P(A) \xrightarrow{\cup} P(A)$.
% \vspace{1cm}
%  \end{enumerate}

% Answers (please circle):\vspace{.3cm}
% \begin{enumerate}
% \item \hspace{.5cm}commutative\hspace{.5cm} associative \hspace{.5cm}
%   identity \hspace{.5cm} inverses \vspace{.6cm}
% \item \hspace{.5cm}commutative\hspace{.5cm} associative \hspace{.5cm}
%   identity \hspace{.5cm} inverses \vspace{.6cm}
% \item \hspace{.5cm}commutative\hspace{.5cm} associative \hspace{.5cm}
%   identity \hspace{.5cm} inverses \vspace{.6cm}
% \end{enumerate}


% \newpage
% %%%%%%%%%%%%%%%%%%%%%%%%%%%%%%%%%%%%%%%%%%%%%%%%%%%%%%%%%%%%%%%%%%%%%%%%%%%%%%%%%%%%%%%%%%%%%%%%%%%%%%%%%

% \ques{20}
% Which of the following relations are reflexive/symmetric/transitive?

% \begin{enumerate}
% \item Let $A$ and $B$ be sets and say that $A \sim B$ if
%   $A \cap B$ is empty.

% \item Let $x$ and $y$ be real numbers and define $x \sim y$ if $x = 1$
%   or $y = 1$.
% \item Let $x$ and $y$ be real numbers and define $x \sim y$ if $xy
%   \in \mathbb{Z}$.

% % \item Let $x$ and $y$ be real numbers and define $x \sim y$ if $x = 1$
% %   or $y = -1$.

% % \item Let $x$ and $y$ be real numbers and define $x \sim y$ if $xy
% %   \not \in \mathbb{Z}$.
% \vspace{1cm}
% \end{enumerate}


% Answers (please circle):\vspace{.3cm}
% \begin{enumerate}
% \item \hspace{.4cm} R \hspace{.4cm} S \hspace{.4cm} T
% \vspace{.6cm}
% \item \hspace{.4cm} R \hspace{.4cm} S \hspace{.4cm} T
% \vspace{.6cm}
% \item \hspace{.4cm} R \hspace{.4cm} S \hspace{.4cm} T
% \vspace{.6cm}
% \end{enumerate}


% Prove using induction that for any integer $n\geq 1$, $$2^0 +2^1 + \cdots +2^{n-1} =2^{n}-1.$$

 % Consider the sequence $\{L_n\}$ defined by $L_0 = 0, L_1
 %  = 1$, and $L_{n+2} = (n+1)L_{n+1} + L_{n}$ for $n \geq 1$.

 %  \begin{enumerate}
 %  \item What is $L_4$?\vspace{20pt}
 %  % \item Prove that for every integer $n \geq 1$,
 %  %   \[
 %  %   L_1 + 3L_3 + 5L_5 +...+ (2n-1)L_{2n-1} =L_{2n}.
 %  %   \]

 %  \item Prove that for every integer $n \geq 1$,
 %    \[
 %    2L_2 + 4L_4 + 6L_6 +...+ (2n)L_{2n} =L_{2n+1}-1.
 %    \]
 %  \end{enumerate}



 % Consider the sequence $\{L_n\}$ defined by $L_0 = 0, L_1
 %  = 1$, and $L_{n+2} = (n+1)L_{n+1} + L_{n}$ for $n \geq 1$.

 %  \begin{enumerate}
 %  \item What is $L_4$?\vspace{20pt}
 %  \item Prove that for every integer $n \geq 1$,
 %    \[
 %    L_1 + 3L_3 + 5L_5 +...+ (2n-1)L_{2n-1} =L_{2n}.
 %    \]
 %  \end{enumerate}


 % Consider the sequence $\{L_n\}$ defined by $L_0 = 0, L_1
 %  = 1$, and $L_{n+2} = L_{n+1} + nL_{n}$ for $n \geq 1$.

 %  \begin{enumerate}
 %  \item What is $L_4$?\vspace{20pt}
 %  \item Prove that for every integer $n \geq 1$,
 %    \[
 %    nL_nL_{n+2} = (n+1)L_{n+1} + (-1)^n.
 %    \]
 %  \end{enumerate}




% \newpage
% %%%%%%%%%%%%%%%%%%%%%%%%%%%%%%%%%%%%%%%%%%%%%%%%%%%%%%%%%%%%%%%%%%%%%%%%%%%%%%%%%%%%%%%%%%%%%%%%%%%%%%%%%
% \ques{10} Let $f\colon A \to B$ be a function. Suppose that $f$ is invertible. 
% Prove that $f$ is injective. (Do not use the ``invertible if and only if bijective'' theorem.)
% % Prove that $f$ is surjective.


\newpage
%%%%%%%%%%%%%%%%%%%%%%%%%%%%%%%%%%%%%%%%%%%%%%%%%%%%%%%%%%%%%%%%%%%%%%%%%%%%%%%%%%%%%%%%%%%%%%%%%%%%%%%%%


\begin{center}
% The next line puts horizontal space in the array
\renewcommand{\arraystretch}{2}
\begin{tabular}{|c|c|c|} \hline
\hspace{1cm}Problem\hspace*{1cm}&\hspace{1cm}Points\hspace*{1cm}&\hspace{1cm}Score\hspace*{1cm}\\
\hline \the\scores Total & \thesumscore & \\ \hline
\end{tabular}
\end{center}

\end{document}